% Consolidated guide-selftest02 from f8f4709b0ffac9da4c1fecbf2ce873002f6da6cc; source fragments preserved.

% === BEGIN guide-selftest02-choices.tex blob 9129e876bd7384dfd948eba676a76b6afac10647 ===
% 第15批：《大学物理学学习指导》自测题2选择2.1.1—2.1.9。
% 印刷49—51页／PDF物理57—59页；逐题按扫描页重述、推导。
\subsection{自测题2：质点动力学和刚体力矩选择}

\sourcetest{2.1.1}{49}{57}
\problemtext 质点在平面内的位矢$\bm r=at^2\bm i+bt^2\bm j$（$a,b$常量），判别路径与运动性质。
\workings
两个坐标分量成固定比例$y/x=b/a$（$a\ne0$），轨迹为直线；
$\bm v=2t(a\bm i+b\bm j)$随时间改变，因此是\textbf{变速直线运动}，选\textbf{B}。
\crosscheck $a=0$时轨迹为$y$轴，同样是直线；不能将位矢有两个分量误判为曲线运动。

\sourcetest{2.1.2}{49}{57}
\problemtext 从加速度方向、平均速率与平均速度、恒定速率和速度变化四条说法中选出正确者。
\workings
恒定加速度不保证运动方向恒定，例如平抛运动；
平均速率是路程除时间，平均速度是位移除时间，两者并不恒相等；
一般变速运动的平均速度不能直接由始末速度算术平均求得。
恒定\emph{速率}仍可改变速度方向，例子为匀速圆周运动，
因此选\textbf{D}。
\crosscheck 区分标量速率与矢量速度。

\sourcetest{2.1.3}{49}{57}
\problemtext 与水平成30°的斜向上力$F$把重为$G$的物块压靠竖直墙。
若无论增大$F$多少均不能使物块向上滑动，求静摩擦系数$\mu$的条件。
\workings
法向力$N=F\cos30^\circ$，
切向向上的力为$F\sin30^\circ$。
物块不向上滑的必要条件
$F\sin30^\circ\le G+\mu F\cos30^\circ$；
令$F$增大而$G$固定，得
\[
\boxed{\mu\ge\tan30^\circ=\frac1{\sqrt3}}.
\]
选\textbf{B}。
\crosscheck 不能把重力视作随$F$增大而同步增加的摩擦能力。

\sourcetest{2.1.4}{49}{57}
\problemtext 质量$m_B=2m_A$的物块A、B以压缩弹簧连接，
在光滑水平桌上从静止释放后，求动能$E_{kA}/E_{kB}$。
\workings
系统初始总动量0，释放后$p_A=-p_B$，
因此
\[
\boxed{\frac{E_{kA}}{E_{kB}}
=\frac{p^2/(2m_A)}{p^2/(2m_B)}
=\frac{m_B}{m_A}=2}.
\]
选\textbf{B}。
\crosscheck 动量大小相同、质量较小的一方动能较大。

\sourcetest{2.1.5}{49}{57}
\problemtext 甲乙体重身高相同，各攀同一无摩擦定滑轮轻绳一端，
均从静止开始；虽然甲相对绳的爬速是乙的两倍，问谁先到顶。
\workings
两端轻绳具有相同张力$T$。
若二人质量均为$m$，对地加速度都为$a=T/m-g$；
同初始位置与初速则对地速度与位移始终相同，
\[
\boxed{y_A(t)=y_B(t)}.
\]
因此\textbf{同时到达}，选\textbf{C}。
\crosscheck “相对绳子爬得快”不代表“对地上升快”，
轻绳在两侧的运动方向相反。

\sourcetest{2.1.6}{49--50}{57--58}
\problemtext 源图2.1.3：斜面弹簧劲度系数$k$，
物块初始位置$a$较自然长度端点$O$靠左$x_1$，
沿坡移动至$b$的位移为$x_2$；求弹簧、物块和地球系统的势能变化。
\workings
初态弹簧压缩量$x_1$；末态弹簧有向形变量$x_2-x_1$。
物块上升$x_2\sin\alpha$，故
\[
\boxed{\Delta U
=\tfrac12 k(x_2-x_1)^2-\tfrac12kx_1^2
+mgx_2\sin\alpha}.
\]
选\textbf{C}。
\crosscheck 计算变化量必须减去原有弹性势能，不可只写末态势能。

\sourcetest{2.1.7}{50}{58}
\problemtext 有转动惯量的定滑轮两端挂$m_1<m_2$，
虽此刻轮子逆时针转动，比较左、右绳张力。
\workings
较重的$m_2$使系统角加速度为顺时针，尽管瞬时角速度可逆时针。
设左张力$T_1$、右张力$T_2$，对轮轴
$(T_2-T_1)R=I\alpha_{\rm clockwise}>0$，故
\[
\boxed{T_2>T_1}.
\]
选\textbf{C（右边大）}。
\crosscheck 角速度方向与角加速度方向不必相同。

\sourcetest{2.1.8}{50}{58}
\problemtext 两个力作用于绕定轴的刚体，
判断“均平行转轴”“均垂直转轴”“合力为零”“对轴合力矩为零”
四种条件与轴向力矩之间的关系。
\workings
均平行转轴时各自轴向力矩均为0，命题(1)真。
均垂直转轴时两轴向力矩\emph{可能}抵消，命题(2)真。
合力为零的\emph{力偶}仍可有力矩，命题(3)假。
合力矩为零不要求合力为零，命题(4)假。
故选\textbf{B（1、2正确）}。
\crosscheck 题问的是对\emph{固定轴}的力矩，不是完整三维力矩向量。

\sourcetest{2.1.9}{50--51}{58--59}
\problemtext 质量$M$、长度$L$的均匀细棒以左端O为转轴，
$J_O=ML^2/3$。质量$m$子弹以垂直于杆的$v$
穿过自由端，穿出后仍同向以$v/2$飞行，
求碰后杆角速度。
\workings
短暂穿透期间系统关于O角动量守恒，子弹损失
$mLv-mL(v/2)=mLv/2$，
\[
\frac13ML^2\omega=\frac12mvL,
\qquad
\boxed{\omega=\frac{3mv}{2ML}}.
\]
选\textbf{B}。
\crosscheck 子弹穿透一般耗散机械能，不能使用机械能守恒替代角动量守恒。
% === END guide-selftest02-choices.tex ===

% === BEGIN guide-selftest02-blanks.tex blob 38fae7a46d4e8d6b8e6a25ca9c29a61d6fba50ef ===
% 自测题2填空2.2.1—2.2.8，印刷51／扫描PDF59页。
\subsection{自测题2：碰撞、旋转与势能填空}

\sourcetest{2.2.1}{51}{59}
\problemtext 质量相等的A、B由竖直轻弹簧相隔，
A在弹簧上端，B在下端且由平台C支撑，初态静止。
突然迅速抽走支持板C，求这一瞬间A与B的加速度大小。
\workings
原图中上方A的重力使\emph{弹簧被压缩}；
平衡时弹簧对A向上的力为$mg$，
对下方B向下的力为$mg$。
刚移去C的瞬间，弹簧尚未改变压缩量，
因而上方A仍满足$mg_{\rm spring}-mg=0$，
下方B同时承受自身重力和弹簧向下的力，得到
\[
\boxed{a_A=0},\qquad \boxed{a_B=2g}.
\]
\crosscheck 这不是\emph{拉伸弹簧}问题！
若将原图的压缩误认成拉伸，会把A、B的加速度颠倒。
方向性判断必须以原图上方A压弹簧、下方B由C承托为准。

\sourcetest{2.2.2}{51}{59}
\problemtext 小球沿半径$R$的竖直光滑圆环无摩擦滑动；
整个圆环绕其竖直直径以角速率$\omega$转动。
使小球从环底离开并能在偏离底部的平衡位置停留，需要多大角速率？
\workings
以小球偏离底部的角度$\theta$表示位置，
旋转参考系的有效势能
$U_{\rm eff}=mgR(1-\cos\theta)
-\tfrac12m\omega^2R^2\sin^2\theta$。
非底部平衡点满足
$\cos\theta=g/(\omega^2R)$，故需要
\[
\boxed{\omega>\sqrt{\frac gR}}.
\]
\crosscheck 等号处仅为底部临界稳定性；
要求真正\emph{偏离底部}的平衡需严格大于。

\sourcetest{2.2.3}{51}{59}
\problemtext 水平面上两小球完全非弹性碰撞并粘合。
$m_1=2.0$g，$\bm v_1=10\bm i$cm/s，
$m_2=5.0$g，$\bm v_2=(3\bm i+5\bm j)$cm/s，
求共同速度大小及与$x$轴的夹角。
\workings
动量矢量守恒：
\[
\bm v=\frac{2(10\bm i)+5(3\bm i+5\bm j)}{7}
=5\bm i+\frac{25}{7}\bm j\quad(\mathrm{cm/s}).
\]
因此
\[
\boxed{v=\sqrt{25+(25/7)^2}\simeq6.15\,\mathrm{cm/s}},
\qquad
\boxed{\alpha=\arctan(5/7)\simeq35.54^\circ}.
\]
\crosscheck 动能不守恒，须按\emph{矢量}动量分别求$x,y$分量。

\sourcetest{2.2.4}{51}{59}
\problemtext 质量$m$小球以$v_0$追上远大于它质量的挡板，
挡板沿同方向以$v<v_0$运动。垂直弹性碰撞后，
求球的速度大小以及板对球的冲量大小。
\workings
极重挡板速度碰撞前后可视为恒定。
在随板运动的惯性系，球的入射相对速度为$v_0-v$，
弹性反弹后相对速度变成$-(v_0-v)$。
回到地面系，球碰后有向速度为
$v_f=v-(v_0-v)=2v-v_0$，
\[
\boxed{|v_f|=|2v-v_0|},\qquad
\boxed{J=m|v_f-v_0|=2m(v_0-v)}.
\]
\crosscheck 不能直接把碰后有向速度$2v-v_0$
当成非负\emph{速率}；当$v_0>2v$小球反向运动。

\sourcetest{2.2.5}{51}{59}
\problemtext 竖直轻弹簧原长时下端小球$m$恰接触地面，
把弹簧上端缓慢提起，直到球刚要脱离地面。
求外力的功。
\workings
球离地临界时弹簧向上拉力$kx=mg$，
所以伸长$x=mg/k$。
之前球未上升，重力势能不变，
弹簧储能
\[
\boxed{W_{\rm ext}=\frac12kx^2=\frac{m^2g^2}{2k}}.
\]
\crosscheck 这是\emph{准静态}提升，忽略最终动能；
不能把此临界状态与放手自由振动最低点混淆。

\sourcetest{2.2.6}{51}{59}
\problemtext 质点质量$m$在吸引型平方反比中心力
$\bm F=-k\bm e_r/r^2$下以半径$r$作圆周运动，
求速度大小及取无穷远零势能时的机械能。
\workings
向心条件$mv^2/r=k/r^2$，得
\[
\boxed{v=\sqrt{\frac{k}{mr}}}.
\]
势能$U=-k/r$，动能$K=k/(2r)$，
所以
\[
\boxed{E=K+U=-\frac{k}{2r}}.
\]
\crosscheck 圆轨道束缚态总能量为负，源题指定无穷远处$U=0$。

\sourcetest{2.2.7}{51}{59}
\problemtext 质量$m$的地球卫星在地面上方高度$2R$绕地球作圆轨道，
地球半径$R$、质量$M$、万有引力常数$G$，
求卫星动能及引力势能。
\workings
轨道半径是$r=R+2R=3R$，
圆轨道动能$K=GMm/(2r)$、势能$U=-GMm/r$，
故
\[
\boxed{K=\frac{GMm}{6R}},\qquad
\boxed{U=-\frac{GMm}{3R}}.
\]
\crosscheck 不可将高度$2R$误作圆轨道半径。

\sourcetest{2.2.8}{51}{59}
\problemtext 一圆盘以恒角加速度转动，
初角速度$20\pi$rad/s，转过60转后角速度$30\pi$rad/s，
求角加速度与经过时间。
\workings
角位移$\Delta\theta=60(2\pi)=120\pi$rad，
$\omega_f^2-\omega_0^2=2\alpha\Delta\theta$，故
\[
\boxed{\alpha=\frac{(30\pi)^2-(20\pi)^2}{240\pi}
=\frac{25\pi}{12}\,\mathrm{rad/s^2}}.
\]
时间为
\[
\boxed{\Delta t=\frac{\omega_f-\omega_0}{\alpha}
=\frac{10\pi}{25\pi/12}=4.80\,\mathrm s}.
\]
\crosscheck 60转要换成120π弧度，不能直接将60代入角位移公式。
% === END guide-selftest02-blanks.tex ===

% === BEGIN guide-selftest02-solutions.tex blob 9f2ab25a8bfa7bb23ec65ca42b4202ee419430f4 ===
% 自测题2：计算题2.3.1—2.3.5、问答2.4.1；
% 《大学物理学学习指导》第5版印刷51—53／PDF59—61页。
\subsection{自测题2：变力功、楔块碰撞、滑块小车与旋转}

\sourcetest{2.3.1}{51--52}{59--60}
\problemtext 非胡克弹簧被外力拉伸$x$所需的力
$F(x)=52.8x+38.4x^2$（SI）。
求从$x_1=0.50$m拉至$x_2=1.00$m外力所做的功；
若弹簧置于光滑水平桌面，一端固定、另一端系2.17kg物块，
从$x_2$静止释放，回到$x_1$时的速度；判断是否保守力。
\workings
外力功为
\[
\begin{aligned}
W&=\int_{0.50}^{1.00}(52.8x+38.4x^2)\,dx\\
&=\left[26.4x^2+12.8x^3\right]_{0.50}^{1.00}\\
&=(26.4+12.8)-(6.6+1.6)
=\boxed{31.0\,\mathrm J}.
\end{aligned}
\]
若忽略摩擦且弹簧静能转为物块动能，
$\tfrac12(2.17)v^2=W$，故
\[
\boxed{v=\sqrt{\frac{62}{2.17}}\simeq5.35\,\mathrm{m/s}}.
\]
由于弹力只依形变量$x$，可定义
$U(x)=\int_0^xF(s)\,ds$，
弹簧回复力为$-\partial U/\partial x$，
因此该理想弹簧力是\textbf{保守力}，
虽不遵守线性胡克定律。
\crosscheck “不满足胡克定律”不等于“不保守”；
若真实弹簧有迟滞及摩擦则另当别论。

\sourcetest{2.3.2}{52}{60}
\problemtext 光滑斜面上质量$m$物块从高$h$静止滑下，
落入与地面有摩擦的沙箱，物块与沙箱合质量$M+m$。
沙箱后方弹簧劲度$k$、初始原长，
碰后最大压缩$l$。
推导地面动摩擦系数$\mu$。
\workings
物块滑至底部前忽略耗散，$v=\sqrt{2gh}$。
水平碰合时动量守恒，沙箱和物块共同速率
\[
V=\frac{m}{M+m}\sqrt{2gh}.
\]
碰后动能为
\[
K=\frac12(M+m)V^2=\frac{m^2gh}{M+m}.
\]
在压缩弹簧的距离$l$内，
机械能消耗于弹簧储能与滑动摩擦：
\[
\frac{m^2gh}{M+m}=\frac12kl^2+\mu(M+m)gl.
\]
解得
\[
\boxed{\mu=
\frac{m^2gh/(M+m)-\tfrac12kl^2}{(M+m)gl}}.
\]
\crosscheck 本式要求右端非负；
如果由给定数据计算为负，则源模型条件
（例如弹簧初态、碰撞过程中的其他功）不自洽。

\sourcetest{2.3.3}{52}{60}
\problemtext 质量$M=10$kg水平小车上有质量$m=1$kg的滑块，
弹簧劲度$k=110\,\mathrm{N/m}$。
初始弹簧压缩$\Delta l=0.05$m、
滑块到右端挡板距离$L=1.1$m；
同时释放且没有摩擦。
求滑块与弹簧刚分离后两物体对地速度及滑块碰上车端的时间。
\workings
总水平动量为零。记滑块向右速度$v$、小车向左速度$-u$，
则$mv=Mu$。
释放的弹簧势能为
\[
\frac12k(\Delta l)^2
=\frac12mv^2+\frac12Mu^2.
\]
代入质量数据，或直接用相对速率
$v_{\rm rel}=v+u$和约化质量
$\mu_r=Mm/(M+m)=10/11$kg，可得
\[
v_{\rm rel}=\Delta l\sqrt{k/\mu_r}
=0.05\sqrt{110(11/10)}=0.55\,\mathrm{m/s}.
\]
两速度由动量守恒确定：
\[
\boxed{v=\frac{M}{M+m}v_{\rm rel}=0.50\,\mathrm{m/s}},
\quad
\boxed{u=\frac{m}{M+m}v_{\rm rel}=0.050\,\mathrm{m/s}}.
\]
分离后在车内的相对速度不变，
到右端挡板所需
\[
\boxed{t=\frac L{v+u}=\frac{1.10}{0.55}=2.00\,\mathrm s}.
\]
\crosscheck 挡板与滑块在同一辆运动小车内，
相遇应除以\emph{相对速率}0.55m/s，
不能只除以对地的滑块速度0.50m/s。

\sourcetest{2.3.4}{52}{60}
\problemtext 可绕同一轴独立转动的圆轮A、B质量分别
$m_A=10$kg、$m_B=20$kg，半径$r_A,r_B$。
从边缘切向拉细绳，直到两轮边缘的切向加速度大小相同，
已知$J_A=m_Ar_A^2/2$、$J_B=m_Br_B^2/2$，
求对应拉力$f_A/f_B$。
\workings
对每个轮子，$f_ir_i=J_i\alpha_i$，
切向加速度$a_t=\alpha_ir_i$相同，
于是
\[
f_i=\frac{J_i}{r_i^2}a_t=\frac12m_i a_t.
\]
因此
\[
\boxed{\frac{f_A}{f_B}=\frac{m_A}{m_B}=\frac12}.
\]
\crosscheck 半径消去是由于题目给定
\emph{同样的边缘切向加速度}，
并且转动惯量按薄圆盘$J=mr^2/2$计算。

\sourcetest{2.3.5}{52--53}{60--61}
\problemtext 均匀球体绕直径经光滑轴匀速旋转，
其半径由$R$对称收缩到$R/2$。
绕该轴的转动惯量$J=2mR^2/5$，
求旋转周期的变化。
\workings
外力对轴净力矩为零，
$J\omega$守恒。质量不变，$J'\!/J=(R/2)^2/R^2=1/4$，
因此$\omega'=4\omega$。
周期$T=2\pi/\omega$，故
\[
\boxed{T'=\frac T4}.
\]
\crosscheck 半径减半后角速度增至4倍，
而不是2倍；能量可由球体收缩的内部做功增加，
不要求转动动能守恒。

\subsection{自测题2：摩擦力—倾角曲线问答}

\sourcetest{2.4.1}{53}{61}
\problemtext 质量$m$木块置于粗糙木板上，
木板与水平面的夹角$\theta$从0连续增至90°，
动、静摩擦系数按题意均取$\mu$。
画出摩擦力大小$f$随$\theta$的变化，
并指出开始滑动时的临界倾角$\theta_0$。
\workings
在未滑动的静止阶段，沿斜面重力分量由静摩擦平衡，
$f=mg\sin\theta$，逐渐增大。
恰开始滑动时
$mg\sin\theta_0=\mu mg\cos\theta_0$，所以
\[
\boxed{\theta_0=\arctan\mu}.
\]
开始滑动后，动摩擦力大小为
$f_k=\mu N=\mu mg\cos\theta$，逐渐下降。
因此所画曲线由以下分段式决定：
\[
\boxed{
f(\theta)=
\begin{cases}
mg\sin\theta,&0\le\theta\le\arctan\mu,\\
\mu mg\cos\theta,&\arctan\mu<\theta\le\pi/2.
\end{cases}}
\]
\crosscheck 若实际静摩擦最大系数$\mu_s$
与动摩擦系数$\mu_k$不同，曲线在临界角会出现向下跳变；
这里严格遵守题设给出的单一$\mu$简化模型。
% === END guide-selftest02-solutions.tex ===
